\documentclass[12pt]{article}

\usepackage[a4paper, total={6in, 8in}]{geometry}
\usepackage{textcomp}

\begin{document}
\setlength{\parindent}{0pt}

\def \t {\textrightarrow}
\def \v {\vspace{1em}}
\def \bi {\begin{itemize}}
\def \ei {\end{itemize}}
\def \s[#1] {\section*{#1}}
\def \ss[#1] {\subsection*{#1}}
\def \sss[#1] {\subsubsection*{#1}}

\s[Italo Svevo]
\ss[Vita Nuova]
Il verismo non è + presente, gia dalla prima opera \t che parte da un rimando biografico e personale \\
Il personaggio ricalca molto Svevo \t Alfonso Nitti \t è un inetto, che è una persona incapace di vivere \t non sa misurarsi col mondo \\
Alfonso non fa altro che questo \t è un borghese ma cerca vita + soddisfacente, quindi si fidanza con Annetta (figlia di un personaggio importante della finanza) \\
Si pero a una consapevolezza, perche lui vive una schermaglia non amorosa ma lavorativa con il cognato \t il meccanismo di zeno si ripete \\
Alfonso cerca di arrampicarsi, e quella coscienza (salva ancora in Pirandello, qua inizia a franare) \t infatti Alfonso non regge il peso della vita, e si suicida \t con pero teatralita \t si prefigura il suo funerale \\
Egli si immagina anna che piange al suo funerale, davanti al carro funebre \t il suo desiderio lo porta a figurarsi la disperazone della compagna \\
La differenza con le altre opere è che lui si suicida \t l'inettitudine è coerente con il personaggio, che è un inetto incapace a vivere e incapace di trovare un senso\\
È incapace di guardare se, ed è proiettato a guardare solo l'altro \t l'inettitudine poggia su questo \\
È già presente un contesto differente

\ss[Senilità]
La seconda opera della trilogia \\
Italo svevo muore in un incidente in macchina, anche giovane \\
È un opera di mezzo \\
Ma quale vecchiaia (Svevo ha 36 anni, è ancora giovane) \t ma viene messa a fuoco la vecchiaia interiore, Emilio Brentani è vecchio dentro, nell'anima \\
Le due parole chiave sono ricordo ed illusione in Leopardi \\
Lui vive nel ricordo e nell'illusione, vive una dimensione interiore di sospensione dettata dall'amore profondo e allusivo di altro per Angiolina \t non ha personalita forte sua, non è donna angelo (nonostante il suo nome, ha i capelli biondi e gli occhi blu) \\
Lei è una prostituta e di facili costumi \t si crea una relazione che illude Emilio, che pensa di avere una relazione con lui \\
Brentani vive di illusioni, vive un Angiolina che è nella sua mente \t si basa su una perfezione illusoria che non esiste \\
Brentani si guarda vivere, immagina un Angiolina che è tutt'altro di quella reale \\
Brentani emilio, ha sorella Amalia \t sono rimasti orfani e vivono condizione di solitudine \t e l'approccio reciproco è quasi morboso, lui è quasi un padre per lei e la protegge, diventando osseisvo nei suoi confronti \\
Per lei lui invece è un fratello ma anche un figlio \\
Si introduce la figura di stefano balli, scultore, uomo di mondo, dandy, erede di Sperelli e Dorian gray \\
Angiolina è attratta da lui, con il quale crea una relazione veloce \\
Anche amalia riceve attenzioni da balli, e scambia questi sentimenti per sentimenti veri provati dal Balli \t anche lei vive l'illusione \\
Emilio trova vita nell'amare, e anche Amalia \\
Partecipano a balletti, che simboleggiano l'arte \\
Amalia poi si lascia morire \t si lascia andare in una malattia dell'anima, un'inettitudine che diventa male dell'anima \\
La malattia interiore la porta ad avere crisi respiratorie e febbre, delira chiamando Stefano, esalando l'etere quindi muore \\
C'è il male d'amore, ma anche un male interiore \t la donna che muore per amore \t è una perdente? ??????????? \\
È l'emblema del disfacimento della figura della donna \\
Amalia è il volto fragile dell'amore \t ma è nobilitande questo sentimento \\
Il balli sta al brentani come angiolina sta ad amalia \\
Si assiste alla merceficazione della donna \\
Del verismo è rimasto \t rimane il mondo di trieste, una città poliglotta e di confine
\end{document}
